\begin{lemma}
\label{lemma-lift-nth-roots}
Let $R$ be a ring. Let $I \subset R$ be a locally nilpotent ideal.
Let $n \geq 1$ be an integer which is invertible in $R/I$. Then
\begin{enumerate}
\item the $n$th power map $1 + I \to 1 + I$, $1 + x \mapsto (1 + x)^n$
is a bijection,
\item a unit of $R$ is a $n$th power if and only if its image in $R/I$
is an $n$th power.
\end{enumerate}
\end{lemma}

\begin{proof}
Let $a \in R$ be a unit whose image in $R/I$ is the same as the image
of $b^n$ with $b \in R$. Then $b$ is a unit
(Lemma \ref{lemma-locally-nilpotent-unit}) and
$ab^{-n} = 1 + x$ for some $x \in I$. Hence $ab^{-n} = c^n$ by
part (1). Thus (2) follows from (1).

\medskip\noindent
Proof of (1). This is true because there is an inverse to the
map $1 + x \mapsto (1 + x)^n$. Namely, we can consider the map
which sends $1 + x$ to
\begin{align*}
(1 + x)^{1/n}
& =
1 + {1/n \choose 1}x +
{1/n \choose 2}x^2 +
{1/n \choose 3}x^3 + \ldots \\
& =
1 + \frac{1}{n} x + \frac{1 - n}{2n^2}x^2 +
\frac{(1 - n)(1 - 2n)}{6n^3}x^3 + \ldots
\end{align*}
as in elementary calculus. This makes sense because the series is finite
as $x^k = 0$ for all $k \gg 0$ and each coefficient
${1/n \choose k} \in \mathbf{Z}[1/n]$ (details omitted; observe that
$n$ is invertible in $R$ by Lemma \ref{lemma-locally-nilpotent-unit}).
\end{proof}